\section{A next-generation corrected-DFPT framework}
\label{sec:nextdfpt}

The proposed method is a correction layer, not a replacement for the mature
machinery of DFPT.  It uses DFPT to obtain phonons, the fully self-consistent
Kohn--Sham perturbation, and its Bloch-state geometry, then spends many-body effort
only on a low-dimensional projected operator.

\begin{figure}[t]
  \centering
  \resizebox{\linewidth}{!}{\begin{tikzpicture}[
  node distance=5.5mm and 4.5mm,
  box/.style={draw,rounded corners=2pt,align=center,font=\small,
              minimum height=8.5mm,text width=2.7cm,inner sep=3pt},
  base/.style={box,fill=softblue,draw=deepblue},
  channel/.style={box,fill=softorange,draw=deeporange,text width=2.45cm},
  decision/.style={box,fill=softgreen,draw=deepgreen,text width=2.8cm},
  arrow/.style={-{Latex[length=2.2mm]},thick,draw=black!70}
]
  \node[base] (dfpt) {DFPT baseline\\$D^{\ks}_{\nu\vq}$};
  \node[base,right=of dfpt] (project) {Localized projection\\$P_cD^{\ks}P_c$};

  \node[channel,right=of project,yshift=18mm] (global) {\texttt{global\_charge}\\global identity};
  \node[channel,right=of project,yshift=6mm] (site) {\texttt{site\_charge}\\relative block shifts};
  \node[channel,right=of project,yshift=-6mm] (internal) {\texttt{internal}\\block-traceless};
  \node[channel,right=of project,yshift=-18mm] (nonlocal) {\texttt{nonlocal}\\inter-block};
  \node[box,fill=black!5,draw=black!55,below=29mm of project]
    (twobody) {Separate future space\\two-body $dU/du,dJ/du$};

  \node[decision,right=12mm of internal,yshift=-6mm] (risk) {Channel weights +\\static/dynamic risk};

  \node[base,right=of risk,yshift=27mm] (safe) {DFPT-safe\\calibration};
  \node[decision,right=of risk,yshift=9mm] (static) {Static\\correction};
  \node[box,fill=softorange,draw=deepred,right=of risk,yshift=-9mm] (dynamic) {Dynamic\\correction};
  \node[box,fill=black!5,draw=black!55,right=of risk,yshift=-27mm] (abstain) {Abstain};

  \node[base,below=43mm of risk] (validate) {Validate fixed-basis vertex,\\then observables};
  \coordinate (merge) at ([xshift=6mm]static.east |- validate.east);

  \draw[arrow] (dfpt) -- (project);
  \foreach \x in {global,site,internal,nonlocal}
    \draw[arrow] (project.east) -- (\x.west);
  \foreach \x in {global,site,internal,nonlocal}
    \draw[arrow] (\x.east) -- (risk.west);
  \draw[arrow,dashed] (project.south) -- (twobody.north);
  \draw[arrow] (risk) -- (safe);
  \draw[arrow] (risk) -- (static);
  \draw[arrow] (risk) -- (dynamic);
  \draw[arrow] (risk) -- (abstain);
  \draw[thick,draw=black!70] (safe.east) -| (merge);
  \draw[thick,draw=black!70] (static.east) -| (merge);
  \draw[thick,draw=black!70] (dynamic.east) -| (merge);
  \draw[thick,draw=black!70] (abstain.east) -| (merge);
  \draw[arrow] (merge) -- (validate.east);
\end{tikzpicture}
}
  \caption{Proposed channel-resolved workflow.  DFPT supplies the baseline
  perturbation.  Projection and operator decomposition expose which physical
  channels a phonon drives.  A risk gate selects uncorrected DFPT, a static
  channel correction, or a dynamic calculation.  Validation is performed first
  on matrix elements and only then on observables.}
  \label{fig:framework}
\end{figure}

\subsection{Corrected perturbation and matrix element}

In a diagonal-channel approximation, define
\begin{equation}
  D^{\mathrm{next}}_{\nu\vq}(\omega)
  =D^{\ks}_{\nu\vq}
  +\sum_a\bigl[\mathcal K_a(\vq,\omega)-1\bigr]
   D^{(a)}_{\nu\vq}.
  \label{eq:dnext}
\end{equation}
The predicted quasiparticle coupling is
\begin{equation}
  g^{\mathrm{next}}_{mn\nu}(\vk,\vq,\omega)
  =\langle\psi_{m,\vk+\vq}|
    D^{\mathrm{next}}_{\nu\vq}(\omega)
    |\psi_{n,\vk}\rangle.
  \label{eq:gnext}
\end{equation}
When all $\mathcal K_a=1$, the method reduces exactly to its DFPT baseline.
Because the correction acts on operators before Bloch-state matrix elements are
taken, it preserves the band, momentum, and phase dependence already present in
DFPT.

\subsection{How each channel is corrected}

\paragraph{Charge channel.}
The implementation must enforce the exact uniform common-potential constraint and
pass a gauge/chemical-potential shift test.  At finite $q$, the UEG relation
\begin{equation}
  \mathcal K_{\rho}(q,0)
  \approx \frac{P(q,0)}{\chi_s(q,0)}
  \label{eq:uegprior}
\end{equation}
is a calibratable prior, not a universal crystal formula.  Crystal local-field
matrices, dimensionality, polar long-range fields, and multi-band screening must
be added explicitly.  A practical first model can fit the density kernel to a
small number of high-level $q$ points while imposing the exact $q=0$ anchor.

\paragraph{Local internal channels.}
For orbital splitting or mixing in a correlated subspace, obtain
$\partial_u\Sigma_{ab}(\omega)$ from finite-displacement DFT+DMFT, a linearized
impurity response, or a validated low-energy model.  If
$\eta_a^{\mathrm{dyn}}\ll1$, retain only $\mathcal K_a(\vq,0)$.  Near orbital,
spin-state, or Mott instabilities, the dynamic kernel and off-diagonal channel
mixing must be kept.

\paragraph{Nonlocal channels.}
When a phonon mainly changes hopping, hybridization, or long-range exchange, a
local impurity self-energy is insufficient.  Targeted GW perturbation theory,
finite-difference quasiparticle calculations, or a nonlocal embedding method can
provide selected matrix elements of $\mathcal K_{ab}$ \citep{Li2019}.  Symmetry
and locality can then interpolate those corrections over equivalent bonds and
wavevectors.

\paragraph{Interaction-modulation channels.}
If a displacement changes $U$, $J$, or another downfolded interaction, compute
those derivatives with constrained response and include the resulting two-body
vertex in the low-energy model.  Constrained DFPT provides a precedent for
separating screening processes during downfolding \citep{Nomura2015}.

\subsection{Three-level decision gate}

The framework chooses the least expensive level justified by diagnostics.

\begin{table}[htbp]
\centering
\caption{Decision levels for the proposed method.}
\label{tab:decisiongate}
\begin{tabularx}{\linewidth}{@{}p{2.6cm}XXp{3.2cm}@{}}
\toprule
Level & Physical diagnostic & Calculation & Required check \\
\midrule
DFPT-safe & Charge dominated; adequate reference bands; weak estimated dynamics
& Standard DFPT matrix elements & Gauge anchor and one representative high-level
control \\
Static correction & Internal or nonlocal weight is appreciable, but
$\eta^{\mathrm{dyn}}$ is small & Channel-projected static
$\mathcal K_a(\vq,0)$ & Held-out modes and momenta \\
Dynamic correction & Rapid vertex variation, soft electronic mode, small
coherence scale, or proximity to an instability & Frequency-dependent
$\mathcal K_{ab}(\vq,\omega)$ & Values across the physical phonon-frequency
window \\
Abstain & Outside the validated channel/basis/material domain & No automatic
prediction & Request a higher-level calculation \\
\bottomrule
\end{tabularx}
\end{table}

\subsection{Uncertainty and invariance requirements}

Every reported $g^{\mathrm{next}}$ should include uncertainty from high-level
training data, operator-basis choice, momentum interpolation, and neglected
channel mixing.  The following tests are mandatory:

\begin{itemize}
  \item invariance to unitary rotations within a declared degenerate orbital
  subspace;
  \item the acoustic sum rule and appropriate translational constraints;
  \item the uniform scalar-potential gauge test;
  \item Hermiticity and time-reversal/symmetry relations for
  $D^{\mathrm{next}}$; and
  \item recovery of uncorrected DFPT with $\mathcal K=I$.
\end{itemize}

\subsection{Why this is a new prediction framework}

Existing beyond-DFPT methods demonstrate how to compute particular corrections,
but they remain expensive when applied to all modes and wavevectors.  The proposed
novelty is the combination of (i) an operator-level physical diagnosis, (ii) a
hierarchical choice of correction theory, (iii) a low-dimensional transferable
kernel, and (iv) explicit uncertainty and abstention.  Its success criterion is
not that it matches a full many-body calculation by construction, but that it
predicts which parts of the DFPT matrix element need correction and reaches a
defined accuracy with far fewer high-level calculations.
